% Imporditud: tools/import_eksperimendid.py
% Allikas: Eesti füüsikaolümpiaadi eksperimendi ülesannete
%   kogu (Rael Kalda, 2023), ülesanne Ü178, lahendus L178.
\setAuthor{Eero Uustalu}
\setRound{lõppvoor}
\setYear{2020}
\setNumber{G E1}
\setDifficulty{5}
\setTopic{staatika}
\setVahendid{niit, 2 koormist (polt ja kaarjas metallist viht), millimeeterpaber, 2 iminapaga näpitsat (laua peale, alla või serva külge kinnitamiseks), žiletiteraga sõrmus niidi lõikamiseks}
\setPunktid{14}
\setAllikas{Eksperimendi kogu Ü178, lahendus lk 141}
\setLahendusseis{toores}

\prob{Masside suhe}
Leidke koormiste masside suhe. Keelatud on kasutada mistahes teostuses kangkaalu põhimõtet.

\hint

\solu
Niidiga saab avaldada vaid tõmbavat jõudu ja see ühtib niidi suunaga. Newtoni I seadus väidab, et paigal saab keha olla vaid kui talle mõjuvate jõudude summa on null. Teeme niidi sisse mingi vahega kaks aasa ja riputame koormised nende külge. Seejärel tõstame koormised niitipidi üles hoides niidi otsi laiali selliselt, et koormisi ühendav niidilõik on horisontaalne. Teeme sellest joonise.

FN1 FN2

FV1 FV2 FT2 $\alpha_2$ $\alpha_1$ FT1 FH2 FH1 mg mg m m

Koormiste raskusjõud on vastavalt m1g ja m2g jõud vastavates niidilõikudes \$\textbackslash{}vec\{F\}\$$N_1$ ja \$\textbackslash{}vec\{F\}\$T 1 ning \$\textbackslash{}vec\{F\}\$$N_2$ ja \$\textbackslash{}vec\{F\}\$T 2

kus Newtoni I seadusest tulenevalt :

\$\textbackslash{}vec\{F\}\$T 1 = $-$\$\textbackslash{}vec\{F\}\$T 2 \$\textbackslash{}vec\{F\}\$$N_1$ + \$\textbackslash{}vec\{F\}\$T 1 + m1g = 0 \$\textbackslash{}vec\{F\}\$$N_2$ + \$\textbackslash{}vec\{F\}\$T 2 + m2g = 0

siit tulenevalt on ka

ning \$\textbackslash{}vec\{F\}\$T 1 = $-$\$\textbackslash{}vec\{F\}\$$H_1$ \$\textbackslash{}vec\{F\}\$T 2 = $-$\$\textbackslash{}vec\{F\}\$$H_2$ \$\textbackslash{}vec\{F\}\$V 1 = $-$m1g \$\textbackslash{}vec\{F\}\$V 2 = $-$m2g

$m_1$ = FV 1 = tan $\alpha_1$ m$^2$ FV 2 tan $\alpha_2$ mkaarjas metallist viht = 1,59.

mpolt
\probend
